\documentclass[../../main.tex]{subfiles}
\begin{document}

{\bf [解答]}

与えられた関数
\begin{align}
  \begin{cases}
    f(x) = 1-\sin x \\
    g(x) = \int_{0}^{x}(x-t)f(t)dt
  \end{cases} \label{eq:1}
\end{align}
を考える．

\cref{eq:1}の第二式を変形して
\begin{align}
  g(x)=x\int_0^xf(t)dt-\int_0^xtf(t)dt
\end{align}
だから一階・2階微分を計算して
\begin{align}
  g'(x)
  &=\int_0^xf(t)dt+xf(x)-tf(t) \\
  &=\int_0^xf(t)dt
\end{align}
および
\begin{align}
  g''(x)&=f(t) \ge 0
\end{align}
である．
したがって$g(x)$は下に凸な関数であるから凸不等式から
\begin{align}
  g(x+y)+g(x-y)\ge 2g(x)
\end{align}
となる．よって題意は示された．
\newpage
\end{document}
